\section{模型假设}

\begin{enumerate}
\item 评价标准假设：所有 Makespan、总额外搬运量和命中率均以题目提供的评估程序在固定配置 \texttt{config.txt} 下的输出为准。
\item 核内调度确定性假设：同一个 Task 交给核内调度器，得到的操作顺序与换入换出位置唯一确定，评估结果可重复。
\item 简化模型假设：用于分析和性能参照的简化模型只考虑计算量、依赖关系、边界搬运量和固定同步延迟，不模拟多核对 DDR 带宽的争用和 L1/UB 容量不足引起的换出。
\item 参数确定性假设：计算图中的操作周期数与张量字节数为确定值，不考虑硬件运行时的波动。
\item 空核允许假设：$k$ 核方案可以只使用其中部分核心，未使用核心的子图序列为空。
\end{enumerate}
